\documentclass[11pt,a4paper]{ctexart}
\usepackage[top=25mm,bottom=25mm,left=24mm,right=24mm]{geometry}
\usepackage{graphicx,booktabs,threeparttable,makecell,multirow,tabularx,array,enumitem,amsmath,amssymb,xcolor,float,fancyhdr,lastpage}
\usepackage[section]{placeins}
\usepackage{etoolbox}
\pretocmd{\subsection}{\FloatBarrier}{}{}
\usepackage[font=small,labelfont=bf,labelsep=quad,skip=4pt]{caption}
\usepackage[colorlinks=true,linkcolor=black,citecolor=black,urlcolor=blue]{hyperref}
\setmainfont{Times New Roman}
\setsansfont{Arial}
\ctexset{section={format=\Large\heiti\raggedright},subsection={format=\large\heiti\raggedright},
         subsubsection={format=\normalsize\heiti\raggedright}}
\newcolumntype{L}{>{\raggedright\arraybackslash}X}
\newcommand{\hlai}[1]{\textbf{#1}}
\renewcommand{\arraystretch}{1.18}
\setlist{nosep,leftmargin=2em}
\renewcommand{\topfraction}{0.9}\renewcommand{\bottomfraction}{0.8}\renewcommand{\textfraction}{0.07}
\renewcommand{\floatpagefraction}{0.8}\setcounter{topnumber}{3}\setcounter{totalnumber}{5}
\setlength{\emergencystretch}{2em}
\pagestyle{fancy}\fancyhf{}
\lhead{\small PUBG FPS AI 新旧模型对比（背景专项）}\rhead{\small 第 \thepage\ 页 / 共 \pageref{LastPage} 页}
\renewcommand{\headrulewidth}{0.4pt}
\newcommand{\figslot}[4]{%
\begin{figure}[htbp]\centering
\IfFileExists{figures/#1.pdf}{\includegraphics[width=#2]{figures/#1.pdf}}{%
\IfFileExists{figures/#1.png}{\includegraphics[width=#2]{figures/#1.png}}{%
\fbox{\parbox[c][4.2cm][c]{#2}{\centering\small\color{gray}图位：\texttt{\detokenize{#1}}\\[2pt]按《图表需求说明》制作后放入 figures 目录}}}}
\caption{#3}\label{fig:#4}\end{figure}}
\newcommand{\keybox}[1]{\par\noindent{\parbox{\linewidth}{#1}}\par}


\definecolor{ai}{HTML}{1F4E79}
\newcommand{\tagbox}[1]{\colorbox{ai!10}{\strut\heiti\color{ai}\ #1\ }}
\newcommand{\greybox}[1]{\par\noindent\colorbox{black!6}{\parbox{\dimexpr\linewidth-2\fboxsep}{\small #1}}\par}
\newcommand{\bignum}[2]{\parbox[t]{0.24\linewidth}{\centering{\fontsize{17}{20}\selectfont\heiti\color{ai}#1}\\[2pt]{\scriptsize #2}}}

\begin{document}
\begin{titlepage}
\centering\vspace*{40mm}
{\Huge\heiti PUBG FPS AI\par}\vspace{6pt}
{\huge\heiti 新旧模型对比\par}
\vspace{8mm}
{\Large 背景声误增强 · 专项简报\par}
\vspace{18mm}
\begin{tabular}{rl}
官方主口径： & 优化背景声的误增强\\
范围： & PUBG 15 段；背景空隙 6 段\\
对照： & 旧 FPS AI vs 新模型（同源原始轨）\\
版本： & v2.0\quad 2026 年 10 月\\
代码仓库： & \url{https://github.com/21377241/fosi-c3-test}\\
\end{tabular}
\vfill
{\small 独立补充；不改动主报告结论。分析代码见上述仓库。\par}
\end{titlepage}

\section{结论}
\keybox{\textbf{支持：真正底噪上的误增强已压到阈值以下}\\[4pt]
以段内安静帧（短时电平 $<-65$\,dBFS）衡量底噪：误增强率 17\%$\rightarrow$0\%；
中位增益 $-0.1$$\rightarrow$$-3.5$\,dB。\\[2pt]
\textbf{其他能力：}主能力（脚步增强/枪声保持）未明显削弱；事件 FER 与枪声音色有回退}

\begin{center}
\bignum{17\%$\rightarrow$0\%}{底噪误增强率}\\%
\bignum{$-0.1$$\rightarrow$$-3.5$}{底噪中位增益（dB）}\\%
\bignum{$+11.8$$\rightarrow$$+12.3$}{脚步中位（dB）}\\%
\bignum{7\%$\rightarrow$11\%}{事件 FER}
\end{center}

\vspace{2mm}
\begin{table}[htbp]\centering
\caption{主口径（安静帧）}\label{tab:nm_bg_summary}
{\small\begin{threeparttable}
\begin{tabularx}{\linewidth}{Lccc}\toprule
指标 & 旧模型 & 新模型 & $\Delta$ \\
\midrule
底噪误增强率（安静帧 $>$3 dB） & 17\% & 0\% & $-$17\% \\
底噪中位增益（dB） & $-$0.1 & $-$3.5 & $-$3.4 \\
\bottomrule\end{tabularx}
\begin{tablenotes}\footnotesize\item 仅 PUBG，背景空隙 n=6；安静帧 <-65 dBFS。\item 与主报告「背景段误增强率 30\%」（VALORANT+PUBG、整段 RMS）不是同一口径。\end{tablenotes}\end{threeparttable}}
\end{table}


\section{证据：两处误增强样例}
旧模型有两段背景空隙增益超过 3\,dB（test4、test14）。下图对每一段：\textbf{上排}原始 / 旧模型 / 新模型横排（同一行同纵轴，整段振幅一眼可比）；\textbf{下排}由蓝框安静区分别引出放大镜，细看底噪。

\figslot{fig_nm_bg_examples}{0.98\linewidth}{上排横排整段对比；下排各轨放大镜。样例 A（test4）整段误增强已消除；样例 B（test14）含突发声，放大镜内底噪已压下}{nmbgex}

\figslot{fig_nm_bg_claim}{0.88\linewidth}{汇总：底噪误增强率与中位增益（安静帧口径）}{nmbg}

\greybox{样例 A：整段旧 $+3.8$$\rightarrow$新 $-0.5$\,dB。
样例 B：整段 RMS 仍高（旧 $+13.3$ / 新 $+14.0$），因夹带脚步类突发声；放大镜安静段旧 $+6.1$$\rightarrow$新 $+1.0$\,dB，底噪已压下。}

\section{其他能力有没有明显削弱？}
\begin{table}[htbp]\centering
\caption{PUBG 新模型：其他能力快检（非环境音主口径）}\label{tab:nm_other}
{\small\begin{threeparttable}
\begin{tabularx}{\linewidth}{Lcccc}\toprule
能力 & 旧模型 & 新模型 & 变化 & 是否明显削弱 \\
\midrule
脚步中位增益 & $+$11.8 & $+$12.3 & $+$0.5 & 否 \\
事件 FER（不含脚步非脚步） & 7\% & 11\% & 4\% & 有回退（有限） \\
纯净枪声频谱 MSE & 2.0 & 3.3 & $+$1.4 & 有回退（有限） \\
\bottomrule\end{tabularx}
\begin{tablenotes}\footnotesize\item “明显削弱”指脚步增强失效或枪声被大幅抬高；FER/MSE 变差记为有限回退。\end{tablenotes}\end{threeparttable}}
\end{table}


\greybox{脚步与同电平选择性略升或持平，不属于削弱。事件 FER、枪声频谱 MSE 变差，属于连带代价，幅度有限、样本也偏少，不宜写成“能力全面变差”。}

\section{一句话}
\keybox{新模型在 PUBG 上\textbf{兑现了“少抬底噪”}；脚步增强未削弱。
事件级误增强与枪声音色有一定回退，但不构成主能力全面变差——若验收只盯背景口径，可以过；若同时卡 FER/音色，需另议。}

\vspace{6mm}
{\small 复现：\texttt{python analysis/report\_pubg\_newmodel.py}\par}
\end{document}
